\chapter{Nasal Packing (Epistaxis Management)}

\section*{Overview}
Nasal packing is the placement of an absorbent or inflatable material into the nasal cavity to control epistaxis that does not respond to direct pressure or cautery. It is an acute management technique used by emergency and ear, nose, and throat (ENT) clinicians to achieve hemostasis.

\section*{Alternative Names}
Anterior nasal packing; posterior nasal packing; epistaxis packing; nasal tamponade

\section*{Principle}
Packing applies direct mechanical pressure to the bleeding nasal mucosa, promoting clot formation and tamponade of the vessel. Moisture-retaining materials hydrate the mucosa and prevent crusting, while inflatable balloons provide controllable compression. The packing is left in place for 24 to 72 hours until the underlying bleeding source has sealed.

\section*{History}
Nasal packing has been used for centuries, with early practitioners employing cloth and gauze; modern absorbent sponges and inflatable balloons were developed in the 20th century.

\section*{Why It Is Done}
Ordered for epistaxis that persists after initial measures such as digital pressure, vasoconstrictive sprays, and chemical or electrical cautery fail to control bleeding. It is also used to stabilize brisk posterior epistaxis, to manage bleeding in coagulopathic patients, and to protect the nasal mucosa after sinonasal surgery.

\section*{Is This a Primary or Secondary Test or Procedure}
A primary therapeutic procedure for acute epistaxis when conservative measures fail.

\section*{How It Is Performed}
The nose is examined with a headlight and nasal speculum, and the bleeding site is identified after suctioning. An anterior pack (absorbent sponge or ribbon gauze layered from floor to roof of the cavity) is placed along the nasal floor; for posterior bleeding, a balloon catheter or posterior pack is advanced and positioned behind the posterior choana. The pack is then secured externally with tape, and the patient is observed for rebleeding.

\section*{Preparation}
The patient sits upright with the head tilted forward and applies pressure while the nose is inspected. Topical vasoconstrictor and local anesthetic are sprayed. Coagulation studies and blood pressure are reviewed, particularly in anticoagulated or hypertensive patients.

\section*{Contraindications and Precautions}
Placement is avoided or performed with caution in suspected cerebrospinal fluid leak or skull base fracture, because of the risk of intracranial migration. Severe septal deviation, nasal tumors, and uncorrected coagulopathy may require specialty consultation rather than blind packing.

\section*{Risks}
Rebleeding, local pain, nasal trauma, and dislodgement; posterior packing carries risks of aspiration, hypoxia, and the trigeminovagal reflex, and is associated with a rare risk of toxic shock syndrome.

\section*{Aftercare and Recovery}
The pack is left in place for 24 to 72 hours and removed in the clinic once bleeding has stopped. Patients are advised to avoid nose blowing, heavy lifting, and anticoagulant use when possible, and are often started on saline spray or petroleum jelly to keep the mucosa moist.

\section*{Outcome and Follow-up}
Resolution of active bleeding confirms adequate tamponade. Recurrent or continued bleeding after removal may indicate a persistent vessel, coagulopathy, or posterior source requiring further intervention.

\section*{Expected Results}
No active bleeding after pack removal, with a healed nasal mucosa and no recurrence during the observation period.

\section*{Follow-up}
Nasal endoscopy to identify the bleeding source, and evaluation of coagulation status or blood pressure if rebleeding occurs. Patients with recurrent epistaxis may benefit from cautery, embolization, or ligation of the responsible vessel.

\section*{References}
\begin{enumerate}
  \item MedlinePlus. Nosebleed. U.S. National Library of Medicine; 2025.
  \item Current Medical Diagnosis and Treatment. McGraw-Hill; 2025.
\end{enumerate}

\clearpage
